% =============================================================================
%  capitulos/cap2_fundamentacao.tex
%  CAPÍTULO 2 — FUNDAMENTAÇÃO TEÓRICA
%
%  Objetivo: apresentar e justificar as tecnologias, padrões e conceitos
%  que embasam o projeto. Seja direto — não é uma enciclopédia.
%
%  DICA: Muito do conteúdo do TCC1 pode ser reaproveitado aqui.
%  Para cada tecnologia: o que é, para que serve e POR QUE foi escolhida
%  (comparando com alternativas, quando possível).
% =============================================================================

\chapter{Fundamentação Teórica}
\label{cap:fundamentacao}

TODO: Breve parágrafo introdutório do capítulo, descrevendo o que será
abordado e como o conteúdo se relaciona ao projeto desenvolvido.


% -----------------------------------------------------------------------------
%  2.1 TECNOLOGIAS UTILIZADAS
%  Apresente cada tecnologia principal. Não liste todas as bibliotecas —
%  foque nas que têm peso arquitetural ou decisório no projeto.
% -----------------------------------------------------------------------------
\section{Tecnologias Utilizadas}
\label{sec:tecnologias}


% --- Subseção por tecnologia ---
% Duplique o bloco abaixo para cada tecnologia relevante do seu projeto.

\subsection{TODO: Tecnologia 1 (ex.: React)}
\label{subsec:tech-react}

TODO: Explique o que é, o propósito principal da tecnologia e por que
ela foi escolhida para este projeto. Se houver alternativas consideradas,
mencione brevemente por que foram descartadas. Cite a fonte oficial:
\cite{react2023}.


\subsection{TODO: Tecnologia 2 (ex.: Python / Django / FastAPI)}
\label{subsec:tech-backend}

TODO: Idem.


\subsection{TODO: Tecnologia 3 (ex.: PostgreSQL / MongoDB)}
\label{subsec:tech-banco}

TODO: Idem.


% Adicione mais subseções conforme necessário. Exemplos:
% \subsection{Docker}
% \subsection{Node.js}
% \subsection{TypeScript}


% -----------------------------------------------------------------------------
%  2.2 PADRÕES DE ARQUITETURA
%  Explique os padrões que guiam a organização do sistema.
%  Ex.: MVC, REST, Microsserviços, Event-Driven, Clean Architecture...
% -----------------------------------------------------------------------------
\section{Padrões de Arquitetura}
\label{sec:padroes}


\subsection{TODO: Padrão 1 (ex.: MVC — Model-View-Controller)}
\label{subsec:padrao-mvc}

TODO: Explique o padrão, suas camadas e como ele se aplica especificamente
ao seu projeto. Cite referências quando mencionar conceitos formais
\cite{pressman2016}.


\subsection{TODO: Padrão 2 (ex.: API REST)}
\label{subsec:padrao-rest}

TODO: Idem.


% -----------------------------------------------------------------------------
%  2.3 TRABALHOS CORRELATOS
%  Sistemas ou trabalhos similares que já existem.
%  Para cada correlato: o que é, o que faz bem, o que falta — deixando
%  evidente a contribuição e diferencial do seu trabalho.
% -----------------------------------------------------------------------------
\section{Trabalhos Correlatos}
\label{sec:correlatos}

TODO: Apresente 2 a 4 sistemas similares ao que você está desenvolvendo.
Para cada um, aponte funcionalidades principais e limitações que justificam
a necessidade do seu trabalho.

A \cref{tab:correlatos} apresenta um comparativo entre os sistemas
correlatos e o sistema proposto.

\begin{table}[H]
  \caption{Comparativo entre sistemas correlatos e o sistema proposto}
  \label{tab:correlatos}
  \centering
  \begin{tabularx}{\textwidth}{lXXXX}
    \toprule
    \textbf{Funcionalidade}  & \textbf{Sistema A} & \textbf{Sistema B}
                             & \textbf{Sistema C} & \textbf{Proposto} \\
    \midrule
    TODO: Funcionalidade 1   & Sim    & Não    & Parcial & Sim  \\
    TODO: Funcionalidade 2   & Não    & Sim    & Sim     & Sim  \\
    TODO: Funcionalidade 3   & Sim    & Sim    & Não     & Sim  \\
    TODO: Gratuito / Aberto  & Não    & Sim    & Não     & Sim  \\
    \bottomrule
  \end{tabularx}
  \fonte{Elaborado pelo autor.}
\end{table}
